\input{preamble.tex}

\title{Отчет по лабораторной работе № 10. \\Настройка прав доступа пользователей к ресурсам сети}
\author{Сущенко Алина Николаевна \\ Группа: НПИбд-01-23}
\date{2025}

\begin{document}

\maketitle
\newpage

\tableofcontents

\newpage
\section{Цель работы}
Освоить настройку прав доступа пользователей к ресурсам сети с использованием стандартных и расширенных списков контроля доступа (Access Control Lists, ACL) в Cisco Packet Tracer.

\newpage
\section{Выполнение работы}

\subsection{Настройка прав доступа на маршрутизаторе msk-ansusenko-gw-1}
Перед началом настройки администратор \texttt{admin-ansusenko} был подключён к порту 24 коммутатора \texttt{msk-ansusenko-sw-4} и получил статический IP-адрес 10.128.6.200, шлюз 10.128.6.1 и DNS-сервер 10.128.0.5.

\begin{enumerate}
\item Мы добавили ноутбук\texttt{admin-ansusenko} и подключили его к \texttt{msk-ansusenko-sw-4} (Рис.\ref{01}):

\begin{center}
    \centering
    \includegraphics[width=\textwidth]{../images/01.png}
    \captionof{figure}{Обновленная топология сети.}
    \label{01}
\end{center}

\item Показаны настройки сетевого интерфейса администратора \texttt{admin-ansusenko} (Рис.\ref{02}):
  \begin{itemize}
    \item IPv4-адрес: 10.128.6.200
    \item Маска подсети: 255.255.255.0
    \item Шлюз по умолчанию: 10.128.6.1
    \item DNS-сервер: 10.128.0.5
  \end{itemize}
  Администратор правильно сконфигурирован для работы в сети.

\begin{center}
    \centering
    \includegraphics[width=\textwidth]{../images/02.png}
    \captionof{figure}{Настройки интерфейса FastEthernet0 администратора.}
    \label{02}
\end{center}

\item Показано создание расширенного списка доступа \texttt{servers-out} и добавление правила для доступа к web-серверу (Рис.\ref{03}):
Правило разрешает доступ всем пользователям к web-серверу (10.128.0.2) по порту 80.

\begin{center}
    \centering
    \includegraphics[width=\textwidth]{../images/03.png}
    \captionof{figure}{Создание ACL servers-out и применение на интерфейсе f0/0.3.}
    \label{03}
\end{center}

\item Добавлены правила для доступа администратора по FTP и Telnet (Рис.\ref{04}):

\begin{center}
    \centering
    \includegraphics[width=\textwidth]{../images/04.png}
    \captionof{figure}{Добавление правил доступа для администратора.}
    \label{04}
\end{center}

\item Далее показана команда \texttt{ftp 10.128.0.2} с устройства администратора (Рис.\ref{05}):
  FTP-доступ к web-серверу успешно работает.

\begin{center}
    \centering
    \includegraphics[width=\textwidth]{../images/05.png}
    \captionof{figure}{Успешное подключение к FTP-серверу.}
    \label{05}
\end{center}

\item Показана попытка подключения по FTP с устройства \texttt{dep-ansusenko-1} (Рис.\ref{06}):
  Доступ запрещён, что соответствует заданию - FTP разрешён только администратору.

\begin{center}
    \centering
    \includegraphics[width=\textwidth]{../images/06.png}
    \captionof{figure}{Ошибка подключения к FTP с устройства dep-ansusenko-1.}
    \label{06}
\end{center}

\item Добавлены правила для доступа к файловому серверу (10.128.0.3):
  Внутренней сети разрешён доступ по SMB (порт 445), всем - по FTP. (Рис.\ref{07}):

\begin{center}
    \centering
    \includegraphics[width=\textwidth]{../images/07.png}
    \captionof{figure}{Настройка доступа к файловому серверу.}
    \label{07}
\end{center}

\item Добавлены правила для доступа к почтновому серверу (10.128.0.4) (Рис.\ref{08}):

\begin{center}
    \centering
    \includegraphics[width=\textwidth]{../images/08.png}
    \captionof{figure}{Настройка доступа к почтовому серверу.}
    \label{07}
\end{center}

\item Показана команда \texttt{nslookup www.ansusenko.rudn.ru} и \texttt{ping} по имени. DNS-сервер работает корректно, доступ к web-серверу по имени успешен. (Рис.\ref{09}):

\begin{center}
    \centering
    \includegraphics[width=\textwidth]{../images/09.png}
    \captionof{figure}{Проверка работы DNS-сервера.}
    \label{09}
\end{center}

\item Показано добавление правила для ICMP (Рис.\ref{10}):
  \begin{minted}{cisco}
    msk-ansusenko-gw-1(config-ext-nacl)#1 permit icmp any any
  \end{minted}
  Правило добавлено в начало списка, что демонстрирует управление порядком применения правил.

\begin{center}
    \centering
    \includegraphics[width=\textwidth]{../images/10.png}
    \captionof{figure}{Добавление правила разрешения ICMP.}
    \label{10}
\end{center}

\item Проверяем список всех подключенных tcp (Рис.\ref{11}):

\begin{center}
    \centering
    \includegraphics[width=\textwidth]{../images/11.png}
    \captionof{figure}{Просмотр настроенных списков доступа.}
    \label{11}
\end{center}

\item Показано создание списка \texttt{other-in} и его применение на интерфейсе f0/0.104 на входящий трафик. Разрешены любые действия только для администратора.(Рис.\ref{12}):

\begin{center}
    \centering
    \includegraphics[width=\textwidth]{../images/12.png}
    \captionof{figure}{Настройка доступа для сети Other.}
    \label{12}
\end{center}

\item Показано создание списка \texttt{management-out} для доступа к сети управления оборудованием (10.128.1.0/24)(Рис.\ref{13}):

\begin{center}
    \centering
    \includegraphics[width=\textwidth]{../images/13.png}
    \captionof{figure}{Настройка доступа к сети управления оборудованием.}
    \label{13}
\end{center}

\subsection{Выполнение заданий для самостоятельной работы}

\item Проверили доступ не с администратора (Рис.\ref{14}):
\begin{center}
    \centering
    \includegraphics[width=\textwidth]{../images/14.png}
    \captionof{figure}{Проверка доступа к серверам.}
    \label{14}
\end{center}

\item Теперь доступ с администратора (Рис.\ref{15}):

\begin{center}
    \centering
    \includegraphics[width=\textwidth]{../images/15.png}
    \captionof{figure}{Проверка доступа администратора к серверам.}
    \label{15}
\end{center}

\item Задаем новую топологию в соответствии с заданием самостоятельной работы (Рис.\ref{16}):

\begin{center}
    \centering
    \includegraphics[width=\textwidth]{../images/18.png}
    \captionof{figure}{Добавляем новый ноутбук в топологию сети Павловской.}
    \label{16}
\end{center}

\item Настраиваем права доступа для нового администратора (Рис.\ref{17}):

\begin{center}
    \centering
    \includegraphics[width=\textwidth]{../images/16.png}
    \captionof{figure}{Права доступа.}
    \label{17}
\end{center}

\item Проверяем списки и настройки для нового администратора (Рис.\ref{18}):

\begin{center}
    \centering
    \includegraphics[width=\textwidth]{../images/17.png}
    \captionof{figure}{Проверка доступа администратора к серверам.}
    \label{18}
\end{center}

\item Показаны настройки интерфейса второго администратора (Рис.\ref{19}):
  \begin{itemize}
    \item IPv4-адрес: 10.128.6.201
    \item Маска: 255.255.255.0
    \item Шлюз: 10.128.6.1
    \item DNS: 10.128.0.5
  \end{itemize}

\begin{center}
    \centering
    \includegraphics[width=\textwidth]{../images/19.png}
    \captionof{figure}{Настройки администратора на Павловской.}
    \label{19}
\end{center}

\item Показано добавление правил для администратора 10.128.6.20 и добавление в списки \texttt{other-in} и \texttt{management-out}. (Рис.\ref{20}):

\begin{center}
    \centering
    \includegraphics[width=\textwidth]{../images/20.png}
    \captionof{figure}{Добавление правил для администратора на Павловской.}
    \label{20}
\end{center}

\item На скриншоте (Рис.\ref{21}) показан обновлённый вывод \texttt{show access-lists} с правилами для обоих администраторов (10.128.6.200 и 10.128.6.201).

\begin{center}
    \centering
    \includegraphics[width=\textwidth]{../images/21.png}
    \captionof{figure}{Итоговый вывод ACL после настройки второго администратора.}
    \label{21}
\end{center}

\item С устройства администратора 10.128.6.201 выполнены (Рис.\ref{22}):
  \begin{itemize}
    \item ping 10.128.6.1 - успешен
    \item ping 10.128.0.2 - 75\% успешных ответов
    \item ftp 10.128.0.2 - успешен
    \item ping 10.128.1.1 - успешен
  \end{itemize}
  Второй администратор имеет полный доступ к ресурсам сети.

\begin{center}
    \centering
    \includegraphics[width=\textwidth]{../images/22.png}
    \captionof{figure}{Проверка доступа администратора на Павловской.}
    \label{22}
\end{center}

\end{enumerate}

\newpage
\section{Выводы}
В результате выполнения лабораторной работы приобрели практические навыки по настройке прав доступа пользователей к ресурсам сети с использованием расширенных списков контроля доступа Cisco. Были настроены ACL \texttt{servers-out}, \texttt{other-in} и \texttt{management-out}, обеспечивающие требуемые политики доступа к web-серверу, файловому серверу, почтовому серверу, DNS-серверу, а также ограничения для сети Other и доступ к сети управления оборудованием.

\end{document}
